Time-varying scalar fields describe pressure, fire activity, radiation, pollution, and simulated flow. Their extrema identify intensity regions; merge levels record when these regions connect. Static temporal summaries expose feature strength, lifetime, and connectivity across distant time steps.

Temporal merge tree maps (TMTM) linearize each field through its augmented merge tree and align consecutive profiles through subtree overlap~\cite{Koepp:2023:TMTM}. Spatiotemporal merge tree maps (ST-MTM) introduce a distance-aware one-dimensional skeleton and temporal coherence~\cite{STMTM:2026}. Both generate orders by recursively arranging complete child subtrees. This keeps hierarchical features together, but the resulting axis follows tree connectivity rather than a fixed spatial coordinate.

The distinction matters when position is part of the analysis. The frame-48 Storms case in Supplement S8 exposes this conflict: the source tree groups $A,B$, while $C$ lies between their reference positions along the selected direction. Keeping $A,B$ contiguous excludes their geographic order: at least one feature must move away from its reference position. Admitting that order instead changes the scalar barriers representing their merge relations. The choice concerns both spatial alignment and the pressure threshold at which structures connect.

We ask \textbf{RQ1}: when hierarchy must be preserved, why is reference displacement unavoidable, and how small can it be? \textbf{RQ2}: which hierarchy constraints can be released to improve alignment, and what merge-level change does the resulting map incur? \textbf{Reference-Guided Merge Tree Maps (RG-MTM)} uses a fixed spatial projection as its target. It diagnoses unavoidable displacement before layout, then constructs either a preserving map, RG-MTM-P, or a selectively relaxed map, RG-MTM-R. Both retain the extrema and their values; the relaxed map reports the change in their pairwise merge levels.

Our evaluation also asks \textbf{RQ3}: how do these choices affect reference alignment, scalar connectivity, and motion proxies across temporal, static, and volumetric inputs?

The paper makes three contributions:
\begin{itemize}
\item A positional certificate that separates interval-packing and hierarchy costs, together with conflict triples that identify the obstructing feature relations.
\item An optimal fixed-order barrier construction and exact small-tree position--topology frontiers, connecting a changed leaf order to its minimum scalar cost.
\item A reference-guided construction connecting diagnosis, hierarchy selection, interval layout, and scalar filling, evaluated on five planar temporal sequences, an independent reference-position control, two original-grid volumetric sequences, and static Nucleon and reconstructed Benzene inputs.
\end{itemize}
The evaluation measures reference order, achieved anchor displacement, rendered merge levels, and approach/separation decisions. Source fields, temporal maps, and error records connect these measurements to the feature relations changed by each mode. Sections~\ref{sec:model}--\ref{sec:analysis} establish the model and its guarantees; Section~\ref{sec:method} describes the construction. Section~\ref{sec:eval} tests the mechanisms and shared-input comparisons, Section~\ref{sec:results} examines source-linked results by spatial dimension, and Section~\ref{sec:discussion} discusses interpretation and limits.
